%
% \pdfoutput=1  % Required for arXiv with pdflatex; commented out for tectonic/xetex compilation.

\documentclass[11pt]{article}

\usepackage[]{acl}

\usepackage{times}
\usepackage{latexsym}

\usepackage[T1]{fontenc}

\usepackage[utf8]{inputenc}

\usepackage{microtype}

\usepackage{inconsolata}

\usepackage{graphicx}
\usepackage{booktabs}


\newcommand{\todo}[1]{\textcolor{blue}{\textbf{\small [TODO: #1]}}}


\title{RETUYT-INCO at BEA 2026 Shared Task: \\
Feature-Enriched mDeBERTa for Word Difficulty Prediction}


\author{Santiago Robaina \and {\bf Luis Chiruzzo} \and {\bf Aiala Rosá} \\
Instituto de Computación, Facultad de Ingeniería, Universidad de la República \\
Montevideo, Uruguay}


\begin{document}
\maketitle
\begin{abstract}
We describe the RETUYT-INCO participation in the BEA 2026 Shared Task on Vocabulary Difficulty Prediction for English Learners, a regression task that predicts GLMM psychometric difficulty scores for English target words given an L1 cue (Spanish, German, or Mandarin). We submitted two closed-track systems: a feature-engineered XGBoost regressor for all three L1s, and, for Spanish, a 3-seed ensemble of \texttt{mdeberta-v3-base} fine-tuned with the same handcrafted features prepended as input text tokens. Our best test result is 1.094 RMSE on Spanish (ensemble), a 13.0\% reduction over the XLM-RoBERTa-base closed baseline. Two findings we highlight. First, a LaBSE cross-lingual cosine between the L1 source word and the English target word is the single highest-impact feature in our experiments. Second, feature-only XGBoost, with no neural fine-tuning and no GPU, already beats the XLM-RoBERTa-base closed dev baseline on average across the three L1s (1.273 vs.\ 1.287 RMSE).
\end{abstract}


\section{Introduction}
\label{sec:introduction}

Predicting how hard an English word is for a second-language learner is a core problem for intelligent tutoring and for adaptive reading and vocabulary tools. The BEA 2026 Shared Task on Vocabulary Difficulty Prediction for English Learners~\citep{felice2026bea} casts this as a regression problem: given an English target word and an L1 cue (source-language word, context sentence, and a partial orthographic clue), predict a continuous Generalized Linear Mixed Model (GLMM) psychometric difficulty score. The task covers three L1s (Spanish, German, Mandarin) and two tracks (closed and open). The closed-track baseline is a fine-tuned XLM-RoBERTa-base~\citep{conneau2019xlmroberta}.

The task sits at the intersection of two traditions. On the NLP side, the Complex Word Identification shared tasks~\citep{paetzold-specia-2016-semeval,yimam-etal-2018-report} framed word-level difficulty as classification by non-native readers; SemEval 2021~\citep{shardlow-etal-2021-semeval} replaced this with a continuous Lexical Complexity Prediction target, and BEA 2024 MLSP~\citep{shardlow-etal-2024-bea} extended the setting to multilingual lexical simplification. Handcrafted features very similar to the ones we use here (word length, frequency, morphology, character $n$-grams) have been standard baselines across this line of work, with fine-tuned neural encoders more recently taking the lead. On the psychometric side, vocabulary testing research long predates these benchmarks~\citep{laufer2004testing}, and the Knowledge-based Vocabulary Lists project~\citep{schmitt2021kvl,schmitt2024kvl} provides the response data that underlies the BEA 2026 targets; the GLMM difficulty scores follow the random-item IRT tradition~\citep{deboeck2008rirt,dunn2024rasch}. The most direct predecessor is \citet{skidmore-etal-2025-transformer}, who benchmark transformer architectures for predicting vocabulary-test item difficulty on a KVL-derived corpus; their XLM-RoBERTa setup is adapted as the BEA 2026 closed-track baseline.

We approach the task under a self-imposed lightweight constraint: we target systems that can be trained and served on a single low-end GPU or on CPU, and we restrict ourselves to the closed track. This reflects the compute realities of the Global South, where access to large clusters and API-gated models is limited, and the privacy constraints of on-device educational deployments. We submit two systems for the closed track: (i)~an XGBoost~\citep{chen2016xgboost} regressor over a set of handcrafted string, morphological, POS, frequency, and cross-lingual semantic-similarity features, for all three L1s, and (ii)~for Spanish only, a 3-seed ensemble of \texttt{mdeberta-v3-base}~\citep{he2023debertav3} fine-tuned for regression with the computed features prepended as input text tokens.

Our research question centres on which handcrafted features carry the most signal for cross-lingual vocabulary difficulty prediction. Two findings emerge. First, a LaBSE~\citep{feng2022labse} cross-lingual cosine between the L1 source word and the English target word is the single highest-impact feature we tested, reducing average dev RMSE by 0.091 across ES/DE/CN in our tree-based pipeline (\S\ref{sec:xgb}) and adding a further 0.017 ES dev RMSE when also prepended as an input token to mDeBERTa (\S\ref{sec:mdeberta}). Second, feature-only XGBoost over our full feature set beats the XLM-RoBERTa-base closed dev baseline on average (1.273 vs.\ 1.287), with no neural fine-tuning and no GPU. The official-test picture is mixed: XGBoost beats the test baseline on Mandarin, is within noise on German, and underperforms on Spanish (\S\ref{sec:results}). The mDeBERTa ensemble is reported as our submitted best system, not as an efficiency contribution: it is the same total size as the XLM-RoBERTa baseline per seed, and approximately triples inference cost as an ensemble.


\section{Dataset and Task}
\label{sec:dataset}

The task provides, per L1, 6{,}091 training items and 677 development items, with a held-out test set released at submission time~\citep{felice2026bea}. Each item has an English target word and its part of speech, a partial orthographic clue of the target, an L1 translation, and an L1 context sentence. The target is a continuous GLMM psychometric difficulty score, ranging roughly from $-6$ to $+5$ with lower values indicating harder items. The primary metric is RMSE, with Pearson correlation as a secondary metric. Two tracks are defined: closed (only the provided data and standard NLP resources) and open (external data and LLMs allowed). We participate in the closed track only; the open track was out of scope under our lightweight constraint.

\section{Feature Engineering and XGBoost}
\label{sec:xgb}

\subsection{Feature set}

We extract a set of handcrafted features per item, grouped by type rather than by order of addition. All features are computed on CPU.

\textbf{Target-word features} (English only): word length, one-hot POS over the eight task categories, Zipf word frequency from the \texttt{wordfreq} package~\citep{speer2022wordfreq}, indicators for common English suffixes (\texttt{-tion}, \texttt{-sion}, \texttt{-ing}, \texttt{-ence}, \texttt{-ment}, \texttt{-able}, \texttt{-ness}, \texttt{-ly}) and prefixes (\texttt{dis-}, \texttt{re-}, \texttt{un-}, \texttt{sub-}, \texttt{out-}, \texttt{over-}), maximum consonant cluster length, double-letter indicator, syllable count, and an English latinate-root indicator.

\textbf{Source-word features} (L1 only): source-word length, source-word syllable count, multiword-source indicator, and an accent indicator that matches Spanish and German diacritics (zero for CN).

\textbf{Cross-lingual string similarity} (both words): length difference and ratio, normalized Levenshtein distance, character-bigram Jaccard overlap, Jaro--Winkler similarity~\citep{winkler1990string}, shared prefix and suffix ratio, a first-letter match indicator, and a suffix-transform match. The suffix-transform rules are hand-written per L1 and encode the expected cognate shape (e.g.\ \textit{-tion} $\rightarrow$ \textit{-ción} in ES, 15 rules total; \textit{-tion} $\rightarrow$ \textit{-ierung}, \textit{-ty} $\rightarrow$ \textit{-tät}, \ldots in DE, 18 rules total; empty for CN).

\textbf{Cross-lingual semantic similarity}: \texttt{embed\_cosine}, the cosine similarity between LaBSE~\citep{feng2022labse} sentence embeddings of the English target word and the L1 source word. This is the single feature that most reduces RMSE in our ablations and is the main narrative thread of this paper.

\textbf{Context}: length of the L1 context sentence in characters.

\subsection{Model and protocol}

We fit an XGBoost regressor on the training split and report RMSE on the development split. The progression in Table~\ref{tab:feature_progression} uses linear regression for the first three feature sets and XGBoost for the rest (\texttt{n\_estimators}=300, \texttt{max\_depth}=5, \texttt{learning\_rate}=0.05). The final row (row~7) uses L1-tuned XGBoost parameters from an 80-iteration random search with 5-fold CV on the training split; tuning moved dev RMSE by only $\sim$0.001, so feature engineering dominated.

\begin{table*}[t!]
\centering
\small
\begin{tabular}{llccccc}
\toprule
\textbf{Feature set} & \textbf{Model} & \textbf{ES} & \textbf{DE} & \textbf{CN} & \textbf{Avg} & \textbf{$\Delta$ vs.\ prev} \\
\midrule
1. \texttt{word\_len} & LR & 1.778 & 1.707 & 1.509 & 1.665 & --- \\
2. + edit distance & LR & 1.705 & 1.596 & 1.487 & 1.596 & $-$0.069 \\
3. + POS & LR & 1.663 & 1.536 & 1.493 & 1.564 & $-$0.032 \\
4. + morphology + shared prefix + context & XGBoost & 1.613 & 1.524 & 1.459 & 1.532 & $-$0.032 \\
5. + Jaro--Winkler + $n$-gram + Zipf freq. & XGBoost & 1.464 & 1.421 & 1.257 & 1.381 & $-$0.151 \\
6. \textbf{+ \texttt{embed\_cosine}} (LaBSE) & XGBoost & \textbf{1.350} & \textbf{1.339} & \textbf{1.181} & \textbf{1.290} & \textbf{$-$0.091} \\
7. + L1-specific (v2) & XGBoost\footnotemark & 1.327 & 1.334 & 1.158 & 1.273 & $-$0.017 \\
\midrule
XLM-RoBERTa-base (closed) baseline & XLM-R ($\sim$270M) & 1.357 & 1.328 & 1.175 & 1.287 & --- \\
\bottomrule
\end{tabular}
\caption{Feature progression on dev (RMSE, 3 dp). Rows 1--3 add the target-word length, a first cross-lingual string feature (edit distance), and target POS. Row~4 adds target-word morphology (suffix/prefix indicators, consonant cluster, double-letter) together with the shared-prefix ratio and context length. Row~5 adds further cross-lingual string similarity (Jaro--Winkler, character-bigram overlap, first-letter match) and target-word Zipf frequency. Row~6 (bold) adds the cross-lingual semantic similarity (\texttt{embed\_cosine}). Row~7 adds the L1-specific v2 block and switches to L1-tuned XGBoost hyperparameters. The last row of the table is the task baseline for reference.}
\label{tab:feature_progression}
\end{table*}
\footnotetext{Tuned parameters for row~7: ES uses \texttt{n\_est}=600, \texttt{depth}=3, \texttt{lr}=0.02, \texttt{subsample}=0.85, \texttt{colsample}=0.70, \texttt{alpha}=0.5, \texttt{lambda}=2.0; DE and CN use \texttt{n\_est}=400, \texttt{depth}=3, \texttt{lr}=0.05, \texttt{subsample}=0.85, \texttt{colsample}=0.85, \texttt{alpha}=0.1, \texttt{lambda}=2.0. The row~7 / row~6 delta therefore conflates the v2 feature block with the L1-tuned hyperparameters. The v2 feature block itself is applied to all three L1s, with per-L1 rule sets (15 suffix-transform rules for ES, 18 for DE, empty for CN so the suffix-transform component is inert there; the remaining v2 features apply equally to all three L1s).}

\subsection{Observations}

Table~\ref{tab:feature_progression} walks up our feature set. The largest single-feature addition is \texttt{embed\_cosine}, which drops average dev RMSE by 0.091 (row~5 $\rightarrow$ row~6). The effect is largest on Mandarin (CN dev RMSE 1.257 $\rightarrow$ 1.181, $-$0.076), which is consistent with CN having no meaningful orthographic signal between the L1 and English; LaBSE provides the missing cross-lingual bridge. On ES and DE, \texttt{embed\_cosine} still gives the largest single drop, alongside target-word Zipf frequency and cross-lingual Jaro--Winkler similarity, which exploit cognate structure.

By row~7, feature-only XGBoost reaches 1.273 average dev RMSE, below the XLM-RoBERTa-base closed dev baseline (1.287), on CPU and with no neural fine-tuning. A leave-one-out ablation on the 9 ES-targeted v2 features confirmed each contributes positively (individual deltas 0.0004--0.0057 on 5-fold CV); we omit per-feature details for space. We do \emph{not} report a drop-one ablation of \texttt{embed\_cosine} from the full v2 feature set; the claim that it is the highest-impact feature is based on incremental-addition evidence only.


\section{mDeBERTa with Feature-Enriched Input}
\label{sec:mdeberta}

For Spanish, we fine-tune \texttt{microsoft/mdeberta-v3-base}~\citep{he2023debertav3} for regression (\texttt{num\_labels}=1). The full model is approximately 276M parameters (an 86M transformer backbone plus a $\sim$190M 250K-token word-piece embedding matrix), essentially the same total size as the XLM-RoBERTa-base closed baseline ($\sim$270M).

\textbf{Input format.} We prepend five feature tokens to the encoder input, then concatenate the raw L1 and English fields with \texttt{[SEP]} markers:
\begin{quote}\small\texttt{wlen=N | nedit=N | pos=X | clue=N |\\ esim=N | L1\_word [SEP] L1\_context [SEP]\\ en\_clue [SEP] en\_word}
\end{quote}
The five prepended features (English word length, normalised edit distance, POS, clue length, and \texttt{esim}) are a subset of the XGBoost feature set rendered as plain text. \texttt{esim} is the LaBSE cosine of \S\ref{sec:xgb} (\texttt{embed\_cosine}), now available to the encoder as a numeric input token.

\textbf{Training.} GLMM scores are normalised to zero mean and unit variance at training time and denormalised before RMSE computation. We train with learning rate $2\!\times\!10^{-5}$, a cosine schedule with 10\% warmup, up to 10 epochs with early stopping (patience 3), batch size 32, weight decay 0.01, fp16, and max sequence length 256. Training was performed on a single Google Colab T4 GPU.

\textbf{Ensemble.} We train three seeds ($\{10, 42, 123\}$) and average their predictions. The ensemble triples inference compute and memory: three $\sim$276M-parameter models are loaded in parallel, for approximately 828M parameters and 3$\times$ forward passes per prediction. Each individual seed already beats the XLM-RoBERTa-base baseline on ES dev (Table~\ref{tab:dev_main}); the ensemble is a $\sim$0.02--0.03 dev RMSE refinement on top of that, not the source of the gap to the baseline.

\textbf{Why Spanish only.} We did not extend the mDeBERTa ensemble to all three L1. DE and CN ensemble were skipped for time reasons and our own limitations of knowledge of the languages. Our DE and CN submission is therefore XGBoost only.

\textbf{Test-time retraining.} The submitted ES ensemble is retrained on \texttt{train}+\texttt{dev} combined before predicting the test set. The test RMSE (1.094) is therefore not directly comparable to the dev RMSE of the same configuration (1.103).


\section{Results and Analysis}
\label{sec:results}

\begin{table*}[t!]
\centering
\small
\begin{tabular}{lrccc}
\toprule
\textbf{System} & \textbf{Params (total)} & \textbf{ES} & \textbf{DE} & \textbf{CN} \\
\midrule
XLM-RoBERTa-base (closed baseline) & $\sim$270M & 1.357 & 1.328 & 1.175 \\
\texttt{full\_xgb\_v2\_embed} (our XGBoost, submitted) & tree ensemble (CPU)\footnotemark & 1.327 & 1.334 & 1.158 \\
\addlinespace
mDeBERTa (seed 10) & $\sim$276M & 1.178 & --- & --- \\
mDeBERTa (seed 42) & $\sim$276M & 1.137 & --- & --- \\
mDeBERTa (seed 123) & $\sim$276M & 1.172 & --- & --- \\
\textbf{mDeBERTa 3-seed ensemble} (submitted, ES only) & \textbf{$\sim$828M at inference} & \textbf{1.103} & --- & --- \\
\bottomrule
\end{tabular}
\caption{Development-split RMSE (3 dp). Per-seed Pearson on ES is 0.789 / 0.821 / 0.826; the ensemble reaches 0.827. The mDeBERTa backbone (86M) and the XLM-RoBERTa backbone ($\sim$85M) are essentially the same size; total parameters including the 250K $\times$ 768 word-piece embedding matrix are $\sim$276M and $\sim$270M respectively.}
\label{tab:dev_main}
\end{table*}
\footnotetext{The submitted XGBoost boosters were not persisted; their serialised size on disk was not measured before submission. The model is a gradient-boosted tree ensemble (up to 600 trees of depth 3) trained on $\leq$23 numeric features.}

\begin{table*}[t!]
\centering
\small
\begin{tabular}{llccccc}
\toprule
\textbf{L1} & \textbf{System} & \textbf{RMSE} & \textbf{Pearson} & \textbf{Team rank} & \textbf{$\Delta$ RMSE vs.\ baseline} \\
\midrule
ES & mDeBERTa 3-seed ensemble (submitted) & 1.094 & 0.843 & 9 / 21 & $-$0.163 ($-$13.0\%) \\
ES & XGBoost (secondary submission) & 1.323 & 0.713 & --- & $+$0.066 ($+$5.3\%) \\
DE & XGBoost (submitted) & 1.260 & 0.713 & 18 / 20 & $+$0.002 ($+$0.2\%) \\
CN & XGBoost (submitted) & 1.106 & 0.754 & 18 / 21 & $-$0.034 ($-$3.0\%) \\
\midrule
ES & XLM-RoBERTa-base (closed test baseline) & 1.257 & 0.765 & --- & --- \\
DE & XLM-RoBERTa-base (closed test baseline) & 1.258 & 0.773 & --- & --- \\
CN & XLM-RoBERTa-base (closed test baseline) & 1.140 & 0.753 & --- & --- \\
\bottomrule
\end{tabular}
\caption{Official test-set results, closed track. A negative $\Delta$ RMSE indicates our system is lower than (better than) the baseline. Team rank collapses each team to its best submission per L1; dashes in that column indicate a row shown for comparison only and not counted separately (our ES team rank is determined by the ensemble). The ES XGBoost row is reported to contextualise the dev/test behaviour of the feature-only system.}
\label{tab:test_main}
\end{table*}

\paragraph{On Dev-set.} Feature-only XGBoost beats the XLM-RoBERTa-base closed dev baseline on all three L1s (Table~\ref{tab:dev_main}), using CPU-fit trees and no neural fine-tuning. The submitted ES ensemble reaches 1.103 dev RMSE, a further 0.224 RMSE below our XGBoost row. This gap is not an efficiency result, since per-seed mDeBERTa has essentially the same parameter count as XLM-RoBERTa-base; it is an effect of the feature-enriched input format, target scaling, and 3-seed averaging. Each individual seed already beats the XLM-RoBERTa dev baseline on ES.

\paragraph{On Test-set.} On the official closed-track test leaderboard (Table~\ref{tab:test_main}), our ES ensemble reports 1.094 RMSE / 0.843 Pearson, placing us 9th of 21 participating teams on Spanish (best-per-team collapse), and reducing the test-baseline RMSE by 13.0\%. XGBoost on Mandarin also beats the test baseline (1.106 vs.\ 1.140, $-$3.0\%), placing us 18th of 21 teams. German XGBoost is within noise of the test baseline (1.260 vs.\ 1.258, $+$0.2\%, rank 18/20). ES XGBoost underperforms the test baseline (1.323 vs.\ 1.257, $+$5.3\%), despite beating the dev baseline (1.327 vs.\ 1.357). This dev/test inversion on ES is the clearest sign of distribution shift between the two splits for our feature-based system.

\section{Conclusions}
\label{sec:conclusions}

Among the handcrafted features we tested, the LaBSE cross-lingual cosine between the L1 source word and the English target word carries the most task-relevant signal: adding it reduces average dev RMSE by 0.091 across ES/DE/CN in our XGBoost pipeline, and adds a further 0.017 ES dev RMSE when also included as a prepended input token to mDeBERTa. Feature-only XGBoost, with no GPU, already beats the XLM-RoBERTa-base closed dev baseline on average (1.273 vs.\ 1.287). Our submitted ES system, a 3-seed mDeBERTa ensemble, places 9/21 on the closed-track ES test leaderboard at 1.094 RMSE.

\section{Limitations}
\label{sec:limitations}
The neural ensemble was run for ES only; DE was skipped for compute reasons and CN was skipped by design (Latin-script feature assumptions). The submitted ES ensemble loads $\sim$828M parameters and runs three forward passes per prediction: it is not an efficient system relative to the baseline, and we report it as our best test result rather than as an efficiency contribution. We prioritise breadth over depth and did not run per-L1 mDeBERTa hyperparameter searches. The ES XGBoost shows a clear dev/test inversion that we do not fully explain. We did not explore open-track resources. Finally, the ``highest-impact feature'' claim for \texttt{embed\_cosine} is supported by incremental-addition evidence only; a drop-one ablation from the full v2 feature set would complement the story and is left for future work.

\section*{Acknowledgments}

This paper was funded by \textit{Agencia Nacional de Investigación e Innovación} (ANII, Uruguay), Project No. $FSED\_2\_2023\_1\_179355$.

\bibliography{anthology,custom}


\end{document}
